\section{Related work}
\label{sec:related_work}

\subsection{B-spline parameterization and knot placement}

Classical B-spline fitting alternates or jointly optimizes point parameters,
knot vectors, and control points. Standard references describe basis evaluation,
least-squares fitting, and practical knot construction
\cite{deboor1972,cox1972,piegl1997,dierckx1993}. Parameterization rules such as
chord length and centripetal parameterization remain strong deterministic
baselines \cite{lee1989}. \TODO{Add recent adaptive knot-placement and model
selection papers from Computer-Aided Design and Computer Aided Geometric Design.}

\subsection{Learning geometric representations for CAD}

Recent geometric learning methods infer parametric primitives, CAD structure,
or editable representations from sampled geometry. Large CAD datasets have
enabled supervised learning of geometric quantities and reconstruction
pipelines \cite{koch2019abc}. \TODO{Add the most relevant curve- and
spline-reconstruction methods, clearly distinguishing unordered point-cloud
methods from the ordered-curve setting studied here.}

\subsection{Sparse gates and set prediction}

Hard-Concrete gates provide a differentiable approximation to an $L_0$
penalty \cite{louizos2018}. Set-prediction architectures instead associate
learned queries with targets through assignment losses \cite{carion2020}.
Although these mechanisms are useful for variable-cardinality prediction,
independent gates require probability calibration and a deployment threshold.
Our formulation makes cardinality explicit and conditions continuous knot
prediction on the selected count.
